\section{Evolution in lattice data}

\subsection{General features}

 An exploratory  lattice study
 of the reduced pseudo-ITD  ${\mathfrak M} (\nu, z_3^2)$
 for the valence $u_v-d_v$ parton distribution in the nucleon has  been reported in Ref.
 \cite{Orginos:2017kos}.
 An amazing observation made there was that,
 when plotted as functions of $\nu$, the data both for real
 and imaginary  parts lie close to respective universal curves.
 The data
 show no polynomial $z_3$-dependence for large $z_3$.
 Given that $z_3^2/a^2$ changes in the explored  range
 from 1 to about 200, we  interpret this result as the  {\it  total
 absence}  of higher-twist terms in the reduced
 pseudo-ITD.

 As  explained in Refs.  \cite{Radyushkin:2017cyf,Orginos:2017kos}   and in the Introduction,
 such an outcome corresponds to  a factorization
 of the $\nu$-  and $z_3^2$-dependences of the soft part of the
 Ioffe-time distribution ${\cal  M} (\nu, z_3^2)= M (\nu) {\cal  M} (0, z_3^2)$.
 In terms of TMD  ${\cal F} (x, k_\perp^2)$, this corresponds to factorization
 of its  $x$- and $k_\perp^2$-dependences   in the region of soft   $k_\perp$.
 However, as observed in Ref.  \cite{Orginos:2017kos}, there is quite visible $z_3$-dependence
 for small values of $z_3$, namely, $z_3 \lesssim 6a$, that may be explained by perturbative evolution.


 Let us  consider first the real part.
 It  corresponds  to
  the cosine Fourier transform
  \begin{align}
   {\mathfrak  R} (\nu)  \equiv
{\rm Re} \,  {\mathfrak  M} (\nu)=  & \int_0^1 dx \
\cos (\nu x)  \, q_v  (x)
\label{MC}
     \end{align}
of the function  $q_v(x)$  corresponding to  the valence combination, i.e., the
difference  $q_v(x) = q(x) -\bar q(x)$  of
quark and antiquark   distributions.
 In our case,   $q= u-d$.

    \begin{figure}[tbp]
    \centerline{\includegraphics[width=3.3in]{images/datAll}}
     \caption{Real  part of ${\mathfrak  M} (\nu, z_3^2)$  plotted as a function
     of $\nu =Pz_3$ and
     compared to the curve given by Eqs. (\relax {\ref{MC}}), (\relax {\ref{qV}}).
        \label{realc}}
    \end{figure}




In  Ref. \cite{Orginos:2017kos},  it was  found that the
data for the real part are  very close  (see Fig. \ref{realc})  to the curve
${\mathfrak R}_f ( \nu)$
generated  by    the function
     \begin{align}
    f(x) =  \frac{315}{32} \sqrt{x}  (1-x)^{3}\,  .
    \label{qV}
    \end{align}
This shape  was obtained by   forming  cosine Fourier transforms  of the
normalized $x^a(1-x)^b$-type  functions
and fixing  the parameters $a,b$  through   fitting    the  data.

   \begin{figure}[tbp]
  \centerline{\includegraphics[width=3.3in]{images/data713.pdf}}
     \caption{Real part   of ${\mathfrak  M} (\nu, z_3^2)$  for  $z_3$  ranging  from  $7a$ to  $13a$.
    \label{M713}}
    \end{figure}

\begin{figure}[tbp]
  \centerline{\includegraphics[width=3.3in]{images/data16.pdf}}
     \caption{Real part  of ${\mathfrak  M} (\nu, z_3^2)$ for $z_3$ ranging   from  $a$ to  $6a$.
     \label{M16}}
    \end{figure}



While
 all the data points have been  used in the fit,
  the shape of the curve   is obviously dominated by
the points with  smaller values of  \mbox{ Re\ ${\mathfrak  M} (\nu, z_3^2)$.  }
To give a more detailed illustration, we show in Fig. \ref{M713} the points
corresponding to $z_3$  values in the range $7a \leq z_3 \leq 13a$.
As  one can see,  there  is some scatter for the points
with the largest values of $\nu$ in the region $\nu \gtrsim 10$,
where  the finite-volume effects  become important.
Otherwise,  practically all the points lie on the universal
curve based on $f(x)$.
 In this sense, there is no $z_3$-evolution visible  in the large-$z_3$ data.

In Fig. \ref{M16}, we show the points in the region  \mbox{$a \leq z_3 \leq 6a$}
(note that,   on the lattice,  $z_3=0$ means that also $\nu=0$,
and ${\mathfrak  M} (0, 0)=1$ by definition).
 In this case, all   the points lie
higher than  the universal curve.
 We recall that the perturbative evolution
increases   the real part of the pseudo-ITD
when $z_3$ decreases. Thus,  one may conjecture that
the  observed higher values
of ${\mathfrak R}$ for smaller-$z_3$ points may be a consequence of the evolution.



A  typical pattern of the $z_3$-dependence of the lattice points
is shown in \mbox{Fig. \ref{nu23}}
for a ``magic''  Ioffe-time value $\nu =3\pi/4$ that may be obtained from
five different combinations of $z_3$ and $P$ values
used in Ref. \cite{Orginos:2017kos}.  The shape of the eye-ball fit  line is
given by the incomplete gamma-function $\Gamma (0,z_3^2/30a^2)$.
This function entirely conforms
to the expectation that the  $z_3$-dependence
has   a ``perturbative'' logarithmic $\ln (1/z_3^2)$  behaviour  for small $z_3$,
and
rapidly vanishes  for $z_3$ larger than $6a$.



As expected,
${\mathfrak R} ( \nu, z_3^2)$  decreases    when $z_3$ increases.
We also see that  the evolution ``stops'' for large $z_3$.
In this context, the overall  curve  based on Eq.  (\ref{qV}) corresponds to  the ``low normalization point'',
i.e.,  to   the region, where the perturbative evolution is absent.

 \begin{figure}[tbp]
  \centerline{\includegraphics[width=3.3in]{images/nu3pi.pdf}}
     \caption{Dependence  on $z_3$ for $\nu=3\pi/4\approx2.3562$.
    \label{nu23}}
    \end{figure}

\subsection{Building $\overline{\rm MS}$  ITD}



 Thus, we see   that the data of \mbox{Fig. \ref{nu23}}
 show a logarithmic evolution behavior  in  the small  $z_3$ region.
 Still,    the
 \mbox{$z_3$-behavior }  starts to visibly deviate
from a pure logarithmic $\ln z_3^2$  pattern for $z_3\gtrsim 5a$.
 This  sets the  boundary $z_3 \leq 4a$  on the ``logarithmic region''.
So, let us try to  use Eq. (\ref{MNL}) in that region to construct
the $\overline{\rm MS}$ ITD.


It is instructive to split the contributions in Eq.  (\ref{MNL}), where we will
denote  ${\rm Re} \, {\cal I}  (\nu, \mu^2)  \equiv  {\cal I}_R  (\nu, \mu^2) $.
The first, ``evolution'' part, given by
  \begin{align}
 {\cal I}^{\rm ev} _R  (\nu, \mu^2)    =&  {\mathfrak  R} ( \nu, z_3^2) +  \frac{\alpha_s}{2\pi} \, C_F\
 \int_0^1  dw \,   {\mathfrak  R}  (w \nu,z_3^2)   \nn &
 \times    B(w) \,   \ln \left (z_3^2\mu^2 \frac{  e^{2\gamma_E}}{4} \right )
\
\label{MNLev}
 \end{align}
(recall that  $ {\mathfrak  R} ( \nu, z_3^2) \equiv {\rm Re} \,  {\mathfrak  M} ( \nu, z_3^2$))
corresponds to the leading logarithm approximation used in Ref. \cite{Orginos:2017kos}.
 For $z_3 =2 e^{-\gamma_E}/\mu$, the logarithm vanishes, and  we have
   \begin{align}
 {\cal I}^{\rm ev}_R   (\nu, \mu^2)    =&  {\mathfrak  R} ( \nu, (2 e^{-\gamma_E}/\mu)^2) =
  {\mathfrak  R} ( \nu, (1.12/\mu)^2) \  .
\label{IMev}
 \end{align}


 This  happens, of course, only  if,   for an appropriately chosen $\alpha_s$,
 the \mbox{$ \ln  z_3^2$-dependence}
 of the one-loop correction  cancels  the actual $z_3^2$-dependence of the data,
 visible as scatter in  the data points in Fig. \ref{M16}.
  In Ref. \cite{Orginos:2017kos}, it was found
 that this happens when  $\alpha_s/\pi \approx 0.1$.
 Thus,  Eq. (\ref{MNL})  is accurate  only in the  region,
 where the data show a {\it  logarithmic} dependence on $z_3$, i.e., $z_3 \leq 4a$
 in our case.



\begin{figure}[tbp]
  \centerline{\includegraphics[width=3.4in]{images/bulamulaN.pdf}}
     \caption{Functions $B \otimes  {\mathfrak R}_f$ (upper line) and $L \otimes  {\mathfrak R}_f$
   (lower line)  of Eq. (\ref{MNLo}).
    \label{BMLM}}
    \end{figure}



 Since the difference between
  ${\mathfrak  R}  (w \nu,z_3^2)$  and ${\mathfrak R}_f (w \nu)$
  is ${\cal O} (\alpha_s)$, we may   replace
  ${\mathfrak  R}  (w \nu,z_3^2)$ by    ${\mathfrak R}_f (w\nu)$ in Eq. (\ref{MNLev})
  (recall that ${\mathfrak R}_f (\nu)$
 corresponds  to the  PDF of  Eq. (\ref{qV})).
 The remaining part of   ${\cal I}  (\nu, \mu^2) $
 (where we have already substituted  ${\mathfrak  R}  (w \nu,z_3^2)$  by  ${\mathfrak R}_f (w \nu)$)
    \begin{align}
{\cal I}^{\rm NL}_R   (\nu)    = & \frac{\alpha_s}{2\pi} \, C_F\
 \int_0^1  dw \,   {\mathfrak R}_f (w \nu)   \nn &
 \times  \left  \{   B(w) \
+   \left [4\,   \frac{\ln (1-w)}{1-w} -2 (1-w)  \right ]_+ \right  \}
 \nn & \equiv  \frac{\alpha_s}{2\pi} \, C_F\, \left  [ B \otimes  {\mathfrak R}_f + L \otimes  {\mathfrak R}_f\right ]
\
\label{MNLo}
 \end{align}
is due to  corrections
 beyond the leading logarithm approximation.


 As  we have discussed, the $ L \otimes  {\mathfrak R}_f$  term
 reflects the fact that  the actual scale in the evolution part
 of the vertex diagrams is less than $z_3$.
 To illustrate its impact,
  we show,      in Fig. \ref{BMLM},  the  functions  $B \otimes  {\mathfrak R}_f$ and $L \otimes  {\mathfrak R}_f$.
  One can see that the last one is negative and rather  large.
  Its $\nu$-dependence is similar to that of the $B \otimes  {\mathfrak R}_f$ function.
  In fact,
 in the $\nu <5$ region, we have  $L \otimes  {\mathfrak R}_f \approx -3.5 B \otimes  {\mathfrak  R}_v$.
 Thus, the combined effect of these two terms is close to that of $-2.5 B \otimes  {\mathfrak R}_f$.
 As a result, the inclusion of these terms may be approximately treated as a LLA evolution
 with a  modified rescaling factor. Specifically,
 we may write
    \begin{align}
 {\cal I}_R   (\nu, \mu^2)    \approx &  \, {\mathfrak  R} ( \nu, (2 e^{1.25-\gamma_E}/\mu)^2)
   \approx   \, {\mathfrak  R} ( \nu, (4/\mu)^2)   \  .
\label{IMevL}
 \end{align}
 Thus, the  rescaling factor has  changed by a  factor of 4 compared to the
 original LLA value!



We may use $\mu \approx 4/z_3$  as a guide, but
 the  actual  numerical   calculations should, of course,  be  done  using the ``exact'' Eq. (\ref{MNL}).
To proceed,
we choose the value \mbox{$\mu=1/a$}  which,
 at the  lattice spacing of $0.093$ fm used in Ref. \cite{Orginos:2017kos}
  is approximately  \mbox{2.15 GeV.}
   The  estimate (\ref{IMevL})  tells us  that
  the   ITD  $ {\cal I}_R   (\nu, \mu^2)$    at this scale should be close to the pseudo-ITD
  ${\mathfrak  R} ( \nu, z_3^2) $
   for  \mbox{$z_3\approx 4a$,}  a distance  that is on the border of  the $z_3 \leq 4a$   region.
Taking   the  value
  $\alpha_s/\pi =0.1$ used in Ref. \cite{Orginos:2017kos}  and applying  the full one-loop relation (\ref{MNL})
to the data with  $z_3 \leq 4a$, we generate   the  points for   ${\cal I}_R  (\nu, (1/a)^2) $.

 \begin{figure}[tbp]
  \centerline{\includegraphics[width=3.35in]{images/Re035.pdf}}
     \caption{Function  ${\cal I}_R  (\nu, \mu^2) $  for
     $\mu = 1/a$  calculated using  the data with   $z_3$   from  $a$ to  $4a$.
        The upper curve corresponds to the ITD of the CJ15 global fit PDF.
    \label{Msbar16}}
    \end{figure}





  As seen from   \mbox{Fig. \ref{Msbar16},}
all the  points are close to  some  universal    curve
   with a rather small scatter.
     The  curve itself      was obtained by fitting the points by
 the cosine transform  of a normalized $N x^a (1-x)^b$ distribution,
which gave
   $a=0.35$ and $b=3$.
   The magnitude of the scatter illustrates  the error of the fit for the
   ITD in the $\nu \leq 4$ region.
    In Fig. \ref{Msbar16},  we compare
  our $\mu=1/a$ ITD  with the ITD
 obtained from the global fit PDFs
corresponding to the CJ15   \cite{Accardi:2016qay}  global fit.
 One can see that our ITD is systematically below the curve based on the global fit PDFs.

The reason for the discrepancy may be understood from
Fig. \ref{fabx},
where
 we compare
 the normalized \mbox{$N x^{0.35} (1-x)^3\equiv q_v (x, \mu=2.15$ GeV)}
 distribution to CJ15  \cite{Accardi:2016qay} and MMHT 2014
   \cite{Harland-Lang:2014zoa} global fit PDFs,
 taken at the  scale $\mu=2.15$ GeV.
 Unlike the $\sim  x^{0.35} $ function, these PDFs are singular for small $x$,
 which leads to the enhancement of  ITDs for large and moderate values of $\nu$.

To fit the points for  ${\cal I}_R  (\nu, \mu^2) $, we have used the same
simplest $N x^a (1-x)^b$ Ansatz for the PDF as in Ref. \cite{Orginos:2017kos}.
In principle, one  may use
more complicated models for PDFs  and get  practically the same fitted curve for the ITD
in the $\nu \leq 4$ region,
while a somewhat different curve for PDF $q_v (x)$.
The reason is simple: the  inverse cosine Fourier transform is unique only
when one exactly knows the ITD in the whole $0\leq \nu <\infty$ region.
Performing such a transform from a limited $\nu \leq 4$ region,
one needs to add some assumptions  either
about the behavior of the ITD outside this region or about a functional form of the PDF $q_v (x)$.
We fixed our choice by taking $q_v(x)\sim x^a (1-x)^b$. The   study of how  the
shape of $q_v (x)$  varies if one uses more complicated forms, in particular, those used
in the global fits  \cite{Accardi:2016qay,Harland-Lang:2014zoa}  is an interesting problem
that, however,   goes  beyond the scope of the present paper.


        \begin{figure}[tbp]
  \centerline{\includegraphics[width=3.6in]{images/fx035.pdf}}
     \caption{Curve for $u_v(x)-d_v(x)$  at  $\mu=2.15$ GeV  built from the  data shown in Fig. \ref{Msbar16}
 and   compared to CJ15 and MMHT global fits.
    \label{fabx}}
    \end{figure}

     Comparing to the LLA results
     of Ref. \cite{Orginos:2017kos}, we observe
     that   the large negative one-loop correction  in  Eq. (\ref{MNL})   has  visibly changed
       the extracted PDF, which is  now    further
       from the global fit  PDFs.
       The main reason is that the $z_0=2a$ pseudo-ITD constructed in
       Ref. \cite{Orginos:2017kos}
       was treated there as corresponding to the $\mu\approx 1$ GeV  scale,
       while according to the modified rescaling relation (\ref{IMevL}),
       it   should correspond to $\mu \approx 4$ GeV.
       Hence, to get the $\mu \approx 2$ GeV curve,
       one needs to  evolve it  down in $\mu$.

       Still, the guiding idea of Ref. \cite{Orginos:2017kos},  that the
       $\overline{\rm MS}$ ITDs ${\cal I}_R(\nu,\mu^2)$ can  be obtained from the
       reduced pseudo-ITDs
       ${\mathfrak R} (\nu,z_3^2)$
       by an appropriate rescaling $\mu = k/z_3$,  works with a rather good accuracy
   for all $z_3 \leq 6a$     if one takes $k\approx 4$.
    By this rescaling relation,    the $\mu  =1/2a  \approx 1$ GeV ITD  corresponds
       to the $z_3 \approx  8a$ reduced pseudo-ITD.  As we discussed,
        a boundary point beyond which the evolution stops, is $z_3 \approx 6a$.
     Hence, the  pseudo-ITD   at  this distance  is   given by the
       ITD ${\mathfrak R}_f (\nu)$  corresponding to the
    universal fit  function $f(x)$ of \mbox{Eq.  (\ref{qV}).  }
  This result  may be   also
    obtained  by a direct  numerical calculation based on  \mbox{Eq. (\ref{MNL}). }

     Using Eq. (\ref{MNL}) one may also  evolve the
    $\overline{\rm MS}$   ITD below $\mu=1/2a$, and the  resulting functions  will be
    changing with $\mu$.  On the other hand, the pseudo-ITDs do not change
    with $z_3$ when $z_3 \gtrsim 6a$. Hence,  the
    rescaling connection ${\cal I}_R (\nu, \mu^2) \approx {\mathfrak R} (\nu, (4/\mu)^2)$
    in this region becomes less and less accurate when $\mu$ decreases,
    and eventually makes no sense.




\subsection{Imaginary part}


\begin{figure}[tbp]
  \centerline{\includegraphics[width=3.4in]{images/im713.pdf}}
     \caption{Imaginary  part   of ${\mathfrak  M} (\nu, z_3^2)$  for  $z_3$  ranging  from  $7a$ to  $13a$.
    The curve corresponds to $q (x)  +  \bar q (x)= f(x) + 2 \bar q (x)$, with $f(x)$ given by Eq. (\ref{qV})
    and $\bar q (x)$ given by Eq. (\ref{barq})
    \label{imor713}}
    \end{figure}

Imaginary part of the pseudo-ITD may be considered in a similar way.
 It  corresponds  to
  the sine Fourier transform
  \begin{align}
   % {\mathfrak  J} (\nu)  \equiv
{\rm Im} \,  {\mathfrak  M} (\nu)=  & \int_0^1 dx \
\sin  (\nu x)  \, [q (x)  +  \bar q (x)]
\label{MS}
     \end{align}
of the function   given by the sum   $q(x) + \bar q (x) $    of
quark and antiquark   distributions.  This function
differs from   the valence combination  $q_v(x) = q(x) -\bar q(x)$   by $2 \bar q (x) = 2[ \bar u (x) - \bar d(x)]$.
In Fig. \ref{imor713}, we show the data for large $z_3$ values $z_3 \geq 7a$.
Just like in the case of the real part (see Fig. \ref{M713}), the points with $\nu \lesssim 10$
are close to a universal curve. Representing  $q (x)  +  \bar q (x) =q_v(x) +2 \bar q (x)$
and taking $f(x)$ of Eq. (\ref{qV})   as $q_v (x)$,   we find
  \begin{align}
    \bar q (x) \approx  0.1 \,[20  x \, (1-x)^{3}]\,  .
    \label{barq}
    \end{align}
Note that in  Ref.  \cite{Orginos:2017kos}, the fit was made
for all the $z_3$ points (i.e. the points with $z_3 \leq 6a$ have been also included), and  the overall coefficient
for      $\bar q (x) $  was obtained to be 0.07 rather than 0.1.

In Fig. \ref{imor06}, we show data with $z_3 \leq 4a$. As one can see,
all these  points are below the curve  obtained by fitting the $z_3 \geq 7a$ data.
This is in agreement with the fact
that, in the region $\nu \lesssim 6$, the perturbative evolution
decreases   the imaginary part of the pseudo-ITD
when $z_3$ decreases.
Note that the 1-loop  relation holds for the whole function ${\mathfrak M}= {\rm Re } \, {\mathfrak M} + i \, {\rm Im } \, {\mathfrak M}$.
So, we should just separate there real and imaginary parts, and the construction of
the $\overline {\rm MS}$  function ${\rm Im}\,  {\cal I} (\nu, \mu^2) \equiv {\cal I}_I (\nu, \mu^2)$
proceeds in the same way  as for the real part.

 The results are shown in \mbox{Fig. \ref{imev}.}
Again, all the  points are rather close to a   universal    curve
   with a rather small scatter.  The curve shown corresponds to the sine Fourier transform
   of the sum of the  valence distribution \mbox{$q_v (x, \mu =1/a) =N x^{0.35} (1-x)^3$} obtained
  from the study of the
   real part,  and the antiquark contribution $2 \bar q (x, \mu =1/a$). The latter was found from the fit  to be
   given by $ \bar q (x, \mu =1/a=2.15$ GeV) $= 0.07 [20 x (1-x)^3]$.







\begin{figure}[tbp]
  \centerline{\includegraphics[width=3.5in]{images/imz06.pdf}}
     \caption{Imaginary part   of ${\mathfrak  M} (\nu, z_3^2)$  for  $z_3$  ranging  from  $a$ to  $6a$.
The curve is the same as in Fig. \ref{imor713}.
    \label{imor06}}
    \end{figure}


        \begin{figure}[tbp]
  \centerline{\includegraphics[width=3.4in]{images/imevoS.pdf}}
     \caption{Function  ${\cal I}_I  (\nu, \mu^2) $  for
     $\mu = 1/a$  calculated using  the data with   $z_3$   from  $a$ to  $4a$.
    The curve is described in the text.
    \label{imev}}
    \end{figure}
